% Reused from PoPETS-2027/main.tex lines 241-254 (unchanged except where marked SP).
\section{The Rocq Theorem Prover}\label{coq_theorem}
The Rocq theorem prover~\cite{the_coq_development_team} is an interactive proof assistant that lets users define mathematical objects, state specifications about them, and prove that the definitions satisfy those specifications. Proof writing is the most challenging step; users guide Rocq using \emph{tactics}, a domain-specific language for building proofs incrementally. While Rocq offers some automation (e.g., decision procedures and proof search), human guidance is usually required to complete non-trivial proofs. Its popularity stems from dependent types, which allow encoding arbitrary mathematical statements directly in its logic, and from its ability to extract certified programs to efficient code in OCaml, Haskell, or Scheme.  Rocq has verified many real-world projects, including the CompCert C compiler~\cite{10.1145/1111037.1111042} (used by companies such as AbsInt\footnote{https://www.absint.com/partners.htm}), the Fiat-Crypto~\cite{8835346} library (used in BoringSSL for Chrome, Android, and CloudFlare), and the ConCert smart contract framework~\cite{10.1145/3372885.3373829}.

We use CertiRocq \cite{anand2017certicoq} that turns Rocq programs into WebAssembly (wasm) with machine-checked correctness guarantees. For wasm, CertiRocq extends its extraction pipeline to enable compilation of Rocq developments into \texttt{.wasm} binaries. The wasm backend, originally developed as CertiCoq-Wasm~\cite{meier2025certicoq} and now integrated into CertiRocq, is itself implemented and verified in Rocq with its semantics mechanised against the WebAssembly standard. It takes Rocq as input and compiles it to wasm code by mapping Rocq's primitive integer operations directly to native wasm instructions. The key insight is that programs proven correct within Rocq can be extracted, compiled, and executed in a wasm runtime while preserving their formal semantics. This bridges high-assurance formal verification with practical deployment, allowing certified components such as cryptographic protocols or voting primitives to run securely in web or sandboxed environments without sacrificing the guarantees established in the proof assistant.

\subsection{Reused Components}\label{reused_comp}
Our development builds on existing formal infrastructure where possible: the primality of the Helios parameters is established with the Coqprime library~\cite{thery2007primality}, and the WebAssembly backend is provided by CertiRocq~\cite{anand2017certicoq,meier2025certicoq}. The remaining components---the algebraic hierarchy of groups, fields, and vector spaces, the probability monad (Section~\ref{uniform_dist}), the efficient modular arithmetic for Schnorr groups and prime fields, and all protocol definitions and security proofs---were developed from scratch. This choice was deliberate: our design goal of extracting efficient executable code requires constructive definitions with a small extraction footprint, which precluded the reuse of existing proof-oriented libraries whose definitions are not geared towards computation.

\subsection{Trusted Computing Base}\label{tcb}
It is important to be precise about which components of our pipeline are formally verified and which must be trusted. On the verified side are: all protocol-level definitions together with their machine-checked proofs of completeness, special soundness, and special honest-verifier zero-knowledge; the modular arithmetic for the Schnorr group and prime field, including the Coqprime primality proofs of the 2048-bit and 256-bit Helios parameters; and, on the frontend path, the compilation from Rocq to WebAssembly, since CertiCoq-Wasm is itself implemented and verified in Rocq against the mechanised WebAssembly semantics. The trusted side comprises: the Rocq kernel (as in any Rocq development); the standard extraction mechanism to OCaml, which is not verified, together with the OCaml compiler and runtime; the WebAssembly engine of the browser that executes the extracted voting client; the unverified parser that reads bulletin-board data in our Helios verifier; and the source of randomness---our certified functions take random field elements as explicit inputs, so sampling this randomness happens outside the verified core, and the usual caveats about cryptographic random number generation apply. This trusted computing base is comparable to that of prior extraction-based developments~\cite{287095}, with the notable difference that on the WebAssembly path the unverified extraction step is replaced by a verified compiler, extending the machine-checked guarantees from the proof assistant down to the deployed code. The artefact maps every claim of this paper to its theorem (\texttt{THEOREMS.md}) and ships a script that runs \texttt{Print Assumptions} on all of them: the theorems instantiated at the Helios parameters carry the standard-library axioms for Rocq's primitive integers that the Coqprime certificates compute with (and proof irrelevance in the approval tally instantiation), and three lemmas of the Schulze development that restate completeness over the true margins use functional extensionality (Section~\ref{sec:mech}); every other theorem, including every soundness and zero-knowledge theorem, is closed under the global context.




